\subsection*{先追踪三条坐标轴}
二维的两个基向量变成三维的三个。想象透明方盒从同一角出发的三条棱，分别代表 $e_1,e_2,e_3$。线性变换把盒子变成平行六面体；记录三条棱的新方向，内部每点的位置就由同样的线性组合决定。九个矩阵元素依旧只是三份目的地记录。

绕竖直轴旋转四分之一圈时，横向箭头变为纵向，纵向箭头变为负横向，竖直箭头保持不动。把这三条记录排成列，旋转矩阵便自然出现。应用于 $(2,1,3)$ 时，组合三根新基箭头即可得到 $(-1,2,3)$。

\subsection*{寻找运动中没有改变的量}
如何检查公式真像旋转？可以计算平方长度。变换后是 $(-y)^2+x^2+z^2$，恰等于原来的 $x^2+y^2+z^2$。一次数值算例不能证明一般结论，所以应保留任意的 $x,y,z$。Lean 在这里验证的就是这个一般恒等式。

三维还有一个新现象：不同轴的旋转通常不交换。拿一本书，先绕竖直方向再绕左右方向转动，随后试着倒过来做，会得到不同姿态。不能把二维平面内旋转的交换性推广到所有三维旋转。

本节不要求先掌握一整套旋转参数，只需学会从列读动作。后面遇到体积，回想这个盒子；遇到核，想象它压成面或线；遇到特征向量，寻找变换后仍留在原直线上的方向。三维只是多了一根基，核心推理仍是同一套。
